\chapter{Protokolle}

\section{Einf"uhrung}

Ein Protokoll besteht aus einer Reihe von Schritten, die nacheinander
ausgef"uhrt werden. Es sind mindestens zwei Teilnehmer am Protokoll
beteiligt. Durch das Protokoll soll ein gewisser Zweck erf"ullt
werden. Weiterhin m"ussen noch folgende Bedingungen erf"ullt sein:

\begin{itemize}
  \item Das Protokoll ist jedem Teilnehmer bekannt.
  \item Jeder Teilnehmer ist mit dem Protokoll einverstanden. 
  \item Der Ablauf des Protokolls mu\3 eindeutig sein.
  \item Das Protokoll mu\3 vollst"andig sein (es ist in allen 
        F"allen genau bestimmt, was zu tun ist).
\end{itemize}

Das Ziel eines kryptographischen Protokolls ist es, Abh"oren und
Betr"ugen zu verhindern. Au\3erdem darf kein Teilnehmer mehr erfahren
k"onnen als das, was im Protokoll vorgesehen ist.

\begin{beispiel} Informale Protokolle aus dem t"aglichen Leben:
  \begin{itemize}
    \item Telefonische Bestellung von G"utern;
    \item Bezahlen mit Kreditkarte;
    \item B"urgerentscheid;
    \item Kartenspiele.
  \end{itemize}
\end{beispiel}

Diese Protokolle funktionieren meistens, da man die Teilnehmer kennt,
sieht oder h"ort. Computer brauchen im Gegensatz zu Menschen formale
Protokolle.

%% ************* Protokol mit Notar
\section{Protokoll mit Notar}

Ein {\em Notar} ist eine vertrauensw"urdige Instanz, die keinen der
Teilnehmer bevorzugt. Mit Hilfe einer solchen Instanz lassen sich recht einfach
Protokolle angeben.

\begin{beispiel}[Verkauf eines Autos mit Hilfe eines Notars] 
Der Verkauf kann wie folgt abgewickelt werden:
  \begin{enumerate}
    \item Alice will Bob ein Auto verkaufen.
    \item Bob will mit einem Scheck bezahlen.
    \item Alice gibt den KFZ-Schein dem Notar.
    \item Bob gibt den Scheck an Alice, die ihn dann bei der Bank einl"ost.
    \item Falls der Scheck gedeckt war, gibt der Notar den KFZ-Schein an Bob.
    \item Falls aber der Scheck nicht gedeckt war, mu\3 Alice dies dem Notar 
          beweisen und erh"alt daraufhin den KFZ-Schein vom Notar zur"uck.
  \end{enumerate}
Eine zweite L"osung w"are ein von der Bank zertifizierter Scheck. In
diesem Fall k"onnte man ohne Notar auskommen.
\end{beispiel}

Das obige Beispiel l"a\3t sich auf Rechnernetzwerke "ubertragen. Ein
wesentlicher Nachteil ist es, da\3 bei jeder Transaktion der Notar involviert
ist. Dieser Engpa\3 kann durch den Einsatz mehrerer Notare behoben
werden, was aber die Kosten steigert. Eine andere M"oglichkeit, diesen
Engpa\3 zu beheben, ist es, das Protokoll wie folgt in zwei Teile zu
zerlegen:
\begin{itemize}
  \item{\bf Normalfall}
    \begin{itemize}
      \item Alice und Bob handeln einen Vertrag aus,
      \item Alice unterschreibt,
      \item Bob unterschreibt.
    \end{itemize}
  \item{{\bf Ausnahmefall} (Vertrag wird nicht eingehalten)}
    \begin{itemize}
      \item Alice und Bob erscheinen vor Gericht,
      \item Alice und Bob erl"autern den Fall,
      \item Richter f"allt eine Entscheidung.
    \end{itemize}
\end{itemize}
Dieses Protokoll verhindert nicht einen Betrug, hat aber einen {\em
Abschreckungseffekt}. Das Beste, was man manchmal auch erreichen kann,
ist ein {\em abgesichertes Protokoll}. Dieses Protokoll garantiert,
da\3 keiner der Teilnehmer betr"ugen kann. Sobald sich ein Partner
nicht an das Protokoll h"alt (betr"ugen will), bemerkt das der andere
Teilnehmer und bricht ab.

%% *************** Typen von Angriffen gegen Protokolle
\section{Typen von Angriffen gegen Protokolle}

\begin{itemize}
  \item{\bf passiver Angriff:} \\
    Ein Unbeteiligter h"ort die Nachrichten ab und versucht daraus 
    Information zu gewinnen.
  \item{\bf aktiver Angriff:} \\
    Ein Unbeteiligter "andert die Nachrichten (f"angt beispielsweise 
    Nachrichten ab und schickt andere).
  \item{\bf passiver Betrug:} \\
    Ein Teilnehmer des Protokolls versucht zus"atzliche Informationen 
    zu gewinnen (bei\-spielsweise w"ahlt er bei einer vom Protokoll 
    verlangten Wahl von Zufallszahlen geschickt und nicht zuf"allig).
  \item{\bf aktiver Betrug:} \\
    Ein Teilnehmer h"alt sich nicht an das Protokoll (er gibt 
    beispielsweise einen falschen Namen an).
\end{itemize}

%% ************** Eine erste Public-Key-Idee
\section{Eine erste Public-Key-Idee (Merkles Puzzle)}

Merkles Puzzel basiert auf der Idee, R"atsel zu erzeugen, die f"ur den Sender 
Alice und
den Empf"anger Bob in vertretbarer Zeit zu l"osen sind, nicht aber
f"ur einen Abh"orer Eve. Damit soll das Ziel erreicht werden, da\3 vor
der Kommunikation kein gemeinsamer Schl"ussel zwischen Alice und Bob
ausgetauscht werden mu\3. 

{\bf Vorgehensweise:} 
\begin{enumerate}
  \item Bob erzeugt $2^{20} (\approx 10^6)$ Nachrichten der Form {\bf
    "`R"atselnr.: $x_i$, Schl"ussel: $y_i$"'}, wobei $x_i$ und
    $y_i$ jeweils paarweise verschiedene Zufallszahlen sind. Die $y_i$
    haben eine L"ange von mindestens 40 bit. Bob verschl"usselt diese 
    Nachrichten mit einem symmetrischen Verfahren, wobei er verschiedene 
    Schl"ussel der L"ange 20 bit benutzt, und schickt alles an Alice.
  \item Alice w"ahlt eine der verschl"usselten Nachrichten zuf"allig
    aus und entschl"usselt sie durch einen Brute-force-Angriff
    (Aufwand ungef"ahr $10^6$ Schritte).  Alice erh"alt nach der
    Entschl"usselung die beiden Zufallszahlen $x_0$ und $y_0$.
  \item Alice verschl"usselt die Nachricht $m$ mit dem Schl"ussel
    $y_0$ und schickt das Chiffrat zusammen mit $x_0$ im Klartext an
    Bob.
  \item Bob w"ahlt das zu $x_0$ geh"orende $y_0$ aus seiner Datenbank
    und entschl"usselt das Chiffrat von Alice mit dem Schl"ussel $y_0$.
\end{enumerate}

{\bf Aufwandsbetrachtung: } \\
\begin{tabular}{rcl}
  Bob          & $\approx$ & $2^{20}$ Schritte (Erzeugung der R"atsel) \\
  Senden       & $\approx$ & $2^{20}$ R"atsel von Bob und Chiffrat von Alice\\
  Alice        & $\approx$ & $2^{20}$ Schritte (Brute-force-Angriff) und 
                 Verschl"usselung des Klartextes\\
  Eve          & $\approx$ & $2^{40}$ Schritte (Brute-force-Angriff)
\end{tabular} 

Der Abh"orer Eve mu\3 im schlimmsten Fall $2^{20}$ R"atsel mit
Brute-force entschl"usseln bis er $x_0$ gefunden hat und damit auch
$y_0$. Dazu braucht er aber ungef"ahr $2^{40}$ Schritte, da jedes
R"atsel mit einem 20 bit langen Schl"ussel verschl"usselt wurde.
Es ist f"ur Eve ung"unstiger, das Chiffrat von Alice direkt mit einem
Brute-force-Angriff zu entschl"usseln, da er dazu mehr als
$2^{40}$ Schritte br"auchte, weil die Schl"ussel $y_i$ eine L"ange von mehr als
40 bit haben.

Wenn Alice und Bob 10000 Operationen pro Sekunde ausf"uhren k"onnen, dann
ben"otigen sie jeweils ungef"ahr eine Minute. Eve mu\3 aber bei
vergleichbarer Rechenleistung etwa ein Jahr investieren. Man sieht
also, da\3 der Aufwand f"ur den Abh"orer Eve unvergleichbar schneller
w"achst, wenn die Anzahl der R"atsel und die L"ange der Schl"ussel erh"oht
wird.

%************* Einwegfunktionen
\section{Einwegfunktionen}

Einwegfunktionen bilden einen zentralen Baustein der Public-Key-Protokolle. 

\begin{definition}[Einwegfunktion]
Eine {\em Einwegfunktion} ist eine injektive Funktion $f$, die einfach zu
berechnen ist, deren Umkehrfunktion $f^{-1}$ dagegen sehr schwierig zu
berechnen ist.
\end{definition}

Eine mathematisch exakte Definition w"are: Die Funktion $f$ ist in
polynomialer Zeit zu berechnen, die Umkehrfunktion $f^{-1}$ aber
nicht. Funktionen, die diese Eigenschaft erf"ullen, sind bisher
nicht bekannt.

\begin{definition}[Einwegfunktion mit Fallt"ur] 
Eine {\em Einwegfunktion mit Fallt"ur} ist eine Einwegfunktion $f$,
deren Umkehrfunktion $f^{-1}$ einfach zu berechnen ist, falls man eine
zus"atzliche Information besitzt.
\end{definition}

%************* Hashfunktionen  
\section{Hashfunktionen}

\begin{definition}[Hashfunktion]
Eine {\em Hashfunktion} ist eine Funktion $H$, die eine beliebig lange Eingabe 
auf eine Ausgabe fester L"ange (normalerweise k"urzer) abbildet. 
\end{definition}

\begin{beispiel} 
Das XOR aller Eingabebytes ist eine Hashfunktion.
\end{beispiel}

\begin{definition}[kryptographische Hashfunktion]
  Eine {\em kryptographische Hashfunktion} ist eine Hashfunktion, die
  zus"atzlich noch folgende Bedingungen erf"ullt:
  \begin{itemize}
    \item sie ist einfach zu berechnen,
    \item es ist schwierig zu einer vorgegebenen Ausgabe (Hashwert)
      ein Urbild (eine der h"ochstwahrscheinlich unendlich vielen
      Nachrichten mit diesem Hashwert) zu finden, und 
    \item es ist schwierig zwei Nachrichten mit demselben Hashwert 
      (Kollision) zu finden.
  \end{itemize}
\end{definition}

%% ************************ Hybrid-Kryptosysteme   
\section{Hybrid-Kryptosystem}

Public-Key-Kryptosysteme sind etwa einen Faktor 1000 langsamer als
Blockchiffren bei gleichen Sicherheitsanforderungen. Bei den symmetrischen 
Verfahren mu\3 aber vor der Kommunikation ein geheimer Schl"ussel sicher
ausgetauscht werden. Daher wird in den meisten praktisch eingesetzten
Systemen zun"achst ein geheimer Sitzungsschl"ussel $K$ mit einem
Public-Key-Verfahren "ubermittelt. Danach wird ein schnelleres
symmetrisches Verfahren verwendet, das dann den Sitzungsschl"ussel $K$
benutzt. Diese Systeme hei\3en {\em Hybridsysteme}, da zwei
kryptographische Verfahren eingesetzt werden.
\begin{enumerate}
  \item Bob schickt Alice seinen "offentlichen Schl"ussel $B$.
  \item Alice w"ahlt einen Sitzungsschl"ussel $K$, verschl"usselt ihn
         mit $B$ und schickt $E_B(K)$ an Bob.
  \item Bob entschl"usselt $D_{B'}(E_B(K))=K$ mit seinem geheimen
        Schl"ussel $B'$. Jetzt haben beide den gemeinsamen
        Sitzungsschl"ussel $K$, der nur ihnen bekannt ist.
  \item Ab jetzt l"auft die Kommunikation mit einem schnellen
        symmetrischen Verfahren, das den Sitzungsschl"ussel $K$
        verwendet.
\end{enumerate}

%% ********************* Digitale Signatur
\section{Digitale Signatur}

Eine digitale Signatur soll "ahnlich wie eine handschriftliche Unterschrift 
folgende Eigenschaften erf"ullen:
\begin{itemize}
  \item authentisch, d.h. einer Person eindeutig zugeordnet;
  \item nicht kopierbar, d.h. Teile des Dokuments oder das Dokument sind
    nach der Unterschrift nicht "anderbar;
  \item unabstreitbar;
  \item kurz im Vergleich zum Dokument.
\end{itemize}
Es ist schwierig, elektronisch alle Kritertien gleichzeitig zu realisieren.


\subsection{Signatur mit Hilfe eines Notars und eines symmetrischen Kryptosystems} 

"Ubung: Geben Sie unter Verwendung eines symmetrischen Kryptosystems
und eines Notars ein Signaturverfahren an. Wie werden die folgenden
Aufgaben gel"ost:

\begin{itemize}
   \item Wie "ubermittelt Alice an Bob eine signierte Nachricht?
   \item Wie kann Bob Celia nachweisen, da\3 das Dokument von Alice signiert wurde?
\end{itemize}
Welche Kriterien werden von Ihrem Verfahren erf"ullt?


\subsection{Signatur mit Public-Key-Verfahren}

Public-Key-Verfahren k"onnen auch zur Authentifizierung verwendet
werden, indem man den geheimen Schl"ussel zur Verschl"usselung benutzt.
\begin{enumerate}
  \item Alice verschl"usselt die Nachricht $m$ mit ihrem geheimen
        Schl"ussel $A'$. Sie f"ugt dem Chiffrat $c=D_{A'}(m)$ 
        ihren Namen hinzu und schickt dies an Bob.
  \item Bob entschl"usselt $c$ mit Alices "offentlichen Schl"ussel 
        $A$ und erh"alt die Nachricht $m=E_A(c)$.
\end{enumerate}
Diese Signatur ist kopierbar. Diese Tatsache st"ort nicht bei der 
Unterschrift von Vertr"agen. Werden damit aber Schecks unterschrieben, so 
k"onnen diese mehrfach eingel"ost werden. Es m"ussen Zeitstempel, laufende 
Nummern, etc. eingef"ugt und von der einl"osenden Bank gespeichert werden, um
dieses Problem zu beheben. Diese Unterschrift ist authentisch und
unabstreitbar. Ihre L"ange h"angt von der L"ange der Nachricht $m$ ab.
Sie ist also nicht kurz.

\subsection{Signatur mit Hashfunktionen}

Bei der Signatur mit Public-Key-Verfahren mu\3 die ganze Nachricht $m$
verschl"usselt werden. Dies ist oft mit hohem Rechenaufwand verbunden,
wenn die Nachricht $m$ gro\3 ist. Dieses Problem l"a\3t sich mit Hilfe
einer Hashfunktion beheben.
\begin{enumerate}
  \item Alice berechnet den Hashwert $h=H(m)$ der Nachricht $m$.
  \item Alice verschl"usselt $h$ mit ihrem geheimen Schl"ussel $A'$ 
        und erh"alt ihre Unterschrift $D_{A'}(h)$ der Nachricht $m$.
  \item Alice schickt ihre Nachricht $m$ mit ihrer Unterschrift $D_{A'}(h)$ 
        an Bob.
  \item Bob berechnet den Hashwert der empfangenen Nachricht und
        vergleicht ihn mit $E_{A}(D_{A'}(h))$, wobei $A$ der
        "offentliche Schl"ussel von Alice ist.
\end{enumerate}

{\bf Vorteil:} Deutlich geringerer Rechenaufwand und Trennung von
Nachricht und kurzer Unterschrift. \\
{\bf Nachteil:} Es ist nicht so sicher. Die Sicherheit h"angt von der
Wirksamkeit der Angriffe gegen Hashfunktionen ab.

\subsection{Angriffe gegen Hashfunktionen}

Es soll zu einer Nachricht $m$ eine weitere Nachricht $m'$ gefunden werden, 
so da\3 $H(m)=H(m')$ (Kollision) gilt. Dies kann verwendet werden, um dem
Empf"anger eine falsche Nachricht zuzuspielen, ohne da\3 er dies bemerken kann.
\begin{itemize}
\item{\bf Brute-force-Angriff} \\
   Die Nachricht $m$ und der Hashwert $H(m)$ seien bekannt. Es werden
   solange verschiedene Nachrichten $m'$ ausprobiert, bis $H(m)=H(m')$
   gilt. Durch Ab"andern der Nachricht $m$ an $n$ Stellen (durch
   F"ullworte, Umstellungen) ergeben sich $2^n$ verschiedene Texte
   $m'$. Hat der Hashwert $k$ bit L"ange, so m"ussen etwa $2^k$
   Nachrichten $m'$ ausprobiert werden, bis $H(m)=H(m')$ auftritt.
\item{\bf Meet-in-middle-Angriff} \\
   Man kann Hashfunktionen normalerweise derart modifizieren, da\3 sie
   eine Nachricht von hinten, also vom bekannten Hashwert $H(m)$ der
   Nachricht, r"uckhashen.  Diese Tatsache kann ausgenutzt werden,
   um eine Kollision schneller als mit dem Brute-force-Angriff zu
   finden. Dazu wird die Nachricht $m$ in zwei Teile $m_1$ und $m_2$
   geteilt. Im ersten Teil werden "Anderungen an $k/2$ Stellen
   vorgenommen. Es enstehen dadurch $2^{k/2}$ verschiedene Nachrichten
   $m'_1$ mit Vorw"artshashwerten $H(m'_1)$. Diese Nachrichten werden nach
   dem Hashwert sortiert.

   Mit Hilfe des bekannten Hashwerts $H(m)$ wird solange der  
   R"uckwerts\-hashwert $H_{r"uck}(m'_2,H(m))$ vom entsprechend ver"anderten 
   zweiten Teil $m'_2$ der Nachricht $m$ berechnet, bis eine "Ubereinstimmung
   des R"uckw"artshashwerts $H_{r"uck}(m'_2,H(m))$ mit einem der
   Vorw"artshashwerte $H(m'_1)$ auftritt. Sie kann effizient
   mit Bin"arsuche festgestellt werden.  Aufgrund des 
   Geburtstagsparadoxon sind verh"altnism"a\3ig wenige Ver"anderungen 
   an $m_2$ notwending. Eine Kollisionsnachricht 
   $m'$ ergibt sich durch Hintereinanderschreiben von $m_1'$ und $m_2'$.
\end{itemize}

{\bf Vergleich des Aufwands f"ur k=64}

\begin{itemize}
   \item Brute-force-Angriff \\
     \begin{tabular}{rcl}
       Zeitaufwand     & $\approx$ & $2^{64}$ Schritte ($2^{64}$ $m'$ hashen) \\
       Speicheraufwand & $\approx$ & nur eine Nachricht
     \end{tabular} 
   \item Meet-in-middle-Angriff \\
     \begin{tabular}{rcl}
       vorw"arts hashen      & $\approx$ & $2^{32}$ Schritte ($2^{32}$ $m'_1$ hashen)          \\
       Sortieren der Tabelle & $\approx$ & $32 \times 2^{32}$ Schritte                         \\
       Speicheraufwand       & $\approx$ & $2^{32} \times 8$ byte $=$ 32 GByte f"ur $H(m'_1)$
    \end{tabular}
\end{itemize}
{\bf L"osung des Meet-in-the-middle Problems:} Will man $k$ bit
Sicherheit, so mu\3 der Hashwert $2k$ bit lang sein.


%% *************** Schl"usselaustausch
\section{Protokolle f"ur Schl"usselaustauch} 

Es werden Protokolle ben"otigt, die den sicheren Austausch eines
gemeinsamen Sitzungs\-schl"ussels $K$ zwischen zwei Teilnehmern
koordinieren, um ihnen sp"ater den Einsatz eines schnellen
symmetrischen Verfahrens zu erm"oglichen.

\subsection{Protokoll mit einem symmetrischen Verfahren und einer Schl"usselzentrale (SZ)}\label{SZmitSymmSchl}
  Die Schl"ussel von Alice und Bob sind der SZ bekannt. 
  \begin{enumerate}
     \item Alice teilt der SZ mit, da\3 sie mit Bob kommunizieren 
           will.
     \item Die SZ erzeugt einen Sitzungsschl"ussel $K$, berechnet 
           $E_A(K)$ und $E_B(K)$ und schickt dies an Alice.
     \item Alice berechnet $D_A(E_A(K))=K$.
     \item Alice schickt $E_B(K)$ an Bob.
     \item Bob berechnet $D_B(E_B(K))=K$.
     \item Alice und Bob besitzen jetzt einen gemeinsamen 
           Sitzungschl"ussel $K$, der nur ihnen und der SZ bekannt ist.
  \end{enumerate}
{\bf Problem:} Die SZ kennt die Schl"ussel aller Teilnehmer und alle
Sitzungsschl"ussel und kann somit alle Nachricht abh"oren.

\subsection{Protokoll mit Public-Key und einer Schl"usselzentrale}\label{SZmitPK}
Die "offentlichen Schl"ussel von Alice und Bob sind in einer Datenbank
der Schl"usselzentrale gespeichert.
\begin{enumerate}
   \item Alice holt Bobs "offentlichen Schl"ussel $B$ aus der Datenbank.
   \item Alice erzeugt einen zuf"alligen Sitzungsschl"ussel $K$,
         berechnet $E_B(K)$ und schickt dies mit ihrem Namen versehen
         an Bob.
   \item Bob entschl"usselt $D_{B'}(E_B(K))=K$ mit seinem geheimen 
         Schl"ussel $B'$.
   \item Alice und Bob haben jetzt einen gemeinsamen Sitzungsschl"ussel $K$.
\end{enumerate}
{\bf Vorteil:} Die SZ kennt nur die "offentlichen der Teilnehmer. Ihr sind die
Sitzungschl"ussel nicht bekannt.

Dieses Protokoll bietet noch keine Sicherheit, da sich ein aktiver Angreifer
in einem Man-in-the-middle Angriff zwischen die SZ und die Teilnehmer
schalten kann. Er kann dann den Teilnehmer falsche Schl"ussel schicken,
ohne da"s es von ihnen bemerkt wird.

\subsection{Man-in-the-middle Angriff}
Ein passiver Angreifer Eve kann nur das Public-Key-Verfahren oder das
symmetrische Verfahren angreifen. Eve sollte keine Chance haben, diese
Verfahren zu brechen.  Ein aktiver Angreifer Mallory kann aber die
Leitung unterbrechen und die Kommunikation, die "uber sie l"auft,
kontrollieren.
\begin{enumerate}
  \item Alice schickt ihren "offentlichen Schl"ussel $A$ an
        Bob. Mallory f"angt $A$ ab und schickt seinen "offentlichen
        Schl"ussel $M$ an Bob mit Absender Alice.
  \item Bob schickt seinen "offentlichen Schl"ussel $B$ an
        Alice. Mallory f"angt $B$ ab und schickt seinen "offentlichen
        Schl"ussel $M$ an Alice mit Absender Bob.
  \item Alice verschl"usselt die Nachricht $m$ und schickt das Chiffrat
        $c=(E_M(m))$ an Bob.
  \item Mallory entschl"usselt $D_{M'}(c)=m$, liest die Nachricht
        und leitet $E_B(m)$ weiter an Bob.
  \item Der R"uckweg verl"auft analog.
\end{enumerate}
Die von Alice und Bob ausgetauschten Nachrichten k"onnen von Mallory
nicht nur gelesen, sondern auch in beliebiger Weise manipuliert
werden. Alice und Bob haben keine M"oglichkeit festzustellen, ob
ihre Nachrichten abgeh"ort oder manipuliert wurden.

\subsection{Interlock-Protokoll}
Die besten Public-Key-Verfahren sind also wertlos, wenn die Teilnehmer
falsche Schl"ussel von einem Angreifer Mallory bei einem
Man-in-the-middle-Angriff bekommen. Dieser Schwachpunkt kann
teilweise durch das folgende Protokoll behoben werden.
\begin{enumerate}
  \item Alice schickt Bob ihren "offentlichen Schl"ussel $A$.
  \item Bob schickt Alice seinen "offentlichen Schl"ussel $B$.
  \item Alice berechnet $E_B(m_A)=e_1 e_2$ und schickt die erste
        H"alfte $e_1$ an Bob.
  \item Bob berechnet $E_A(m_B)=f_1 f_2$ und schickt die erste H"alfte
        $f_1$ an Alice.
  \item Alice schickt die zweite H"alfte $e_2$ an Bob.
  \item Bob berechnet $D_{B'}(e_1 e_2)=m_A$. Bob beantwortet {\em sofort}
        die Frage von Alice und schickt die zweite H"alfte $f_2$.
  \item Alice beantwortet Bobs Frage $D_{A'}(f_1 f_2)$.
\end{enumerate}
Ein aktiver Angreifer Mallory kann bis zum 3-ten Schritt alles wie
bisher machen. Dann mu\3 aber Mallory eine neue Nachricht erfinden,
sie verschl"usseln und die erste H"alfte an Bob schicken. Im 4-ten
Schritt mu\3 Mallory nochmals eine neue Nachricht erfinden, sie
verschl"usseln und die erste H"alfte an Alice schicken.  Im 5-ten
Schritt wei\3 Mallory Alices Frage, mu\3 aber die zweite H"alfte der
erfundenen Nachricht an Bob schicken. Im 6-ten Schritt erwartet Alice
eine Antwort von Bob, die jetzt aber Mallory beantworten
mu\3. Mallory wei\3 entweder die Antwort oder mu\3 raten. Der 7-ten
Schritt verl"auft analog. Durch das Interlock-Protokoll wird der
Man-in-the-midde Angriff Mallory erschwert, jedoch nicht unm"oglich
gemacht.

Das Protokoll versagt, wenn  Mallory die beiden Teilnehmer kennt oder wenn
Alice und Bob sich nicht direkt kennen und daher keine geeigneten Fragen
stellen k"onnen, d.h. kein gemeinsames Geheimnis besitzen. Das ist vor
allem dann der Fall, wenn Alice und Bob keine Menschen, sondern
Computer sind, die Fragen nur automatisch erzeugen k"onnen.

%% ******************* Schl"usselaustausch mit Signaturen
\subsection{Schl"usselaustausch mit Signaturen}
Das Interlockprotokoll erschwert einen Man-in-the-middle Angriff,
macht ihn aber nicht unm"oglich. Das Protokoll~\ref{SZmitPK} 
kann mit Hilfe der digitalen Signatur so erweitert werden, da"s
ein Man-in-the-middle Angriff f"ur einen Au"senstehenden Mallory
umm"oglich wird.

Die Schl"usselzentrale unterschreibt die "offentlichen Schl"ussel
von Alice und Bob, wobei sie pers"onlich anwesend sein m"ussen. Diese
von der Schl"usselzentrale signierten Schl"ussel sind in einer
Datenbank zug"anglich. Jeder Teilnehmer kennt den "offentlichen
Schl"ussel der SZ und kann damit ihre Unterschrift pr"ufen. Mallory
kann keinen Man-in-the-middle Angriff durchf"uhren, da Alice
und Bob die Authentizit"at der Schl"ussel "uberpr"ufen k"onnen. 
Mallory kann jetzt nur noch zuh"oren oder die Leitung
st"oren. Die SZ kann nicht wie bei der Vergabe von
Sitzungsschl"usseln im Protokoll~\ref{SZmitSymmSchl} alle Nachrichten
abh"oren, ist also nicht mehr der gro\3e Schwachpunkt im System.

Die Schl"usselzentrale kann jetzt nur noch als einzige Instanz einen
Man-in-the-middle Angriff durchf"uhren.

Mit Hilfe einer SZ lassen sich auch sehr einfach Nachrichten
verschicken. Der "offentliche Schl"ussel von Bob wird von der SZ
geliefert. Alice erzeugt Sitzungsschl"ussel $K$, verschl"usselt
Nachricht $m$ mit symmetrischen Verfahren und schickt
$(E_B(K),E_K(m))$ an Bob. Hier ist wieder wichtig, da\3 der "offentliche
Schl"ussel $B$ signiert ist.

\section{Zero Knowledge Protokolle}

Bei einem Zero-knowledge Protokoll weist ein Beweiser $B$ einem
Verifizierer $V$ nach, die L"osung eines bestimmten Problems $P$ zu
kennen. Bei diesem Nachweis wird keine Information "uber die Struktur
der L"osung preisgegeben. Genauer ist gefordert, da\3 bei polynomial
vielen Abl"aufen des Protokolls nicht mehr Information preisgegeben
wird, als die, die der Verifizierer selbst mit polynomialen Aufwand h"atte
bekommen k"onnen.

Die Zero Knowledge Protokolle k"onnen f"ur Authetifikationsverfahren
benutzt werden.

\begin{beispiel} 
Der Keller hat eine Geheimt"ur. $B$ will $V$ beweisen, da\3 er die
Geheimt"ur "offnen kann, ohne aber $V$ das Geheiminis, wie sie
ge"offnet wird, zu verraten.
  \begin{figure}[htb]
  \centerline{\includeeps{keller.eps}{1.5}}
  \caption{Keller mit Geheimt"ur}
  \end{figure}
  \begin{enumerate}
    \item $V$ steht bei $e$.
    \item $B$ geht entweder zu $g$ oder $h$ im Keller.
    \item $V$ geht zu $f$ und ruft $B$ zu, er solle links bzw. 
          rechts herauskommen.
    \item $B$ "offnet bei Bedarf die Geheimt"ur und kommt am 
          gew"unschten Ort heraus.
    \item Diese Prozedur wird n-mal wiederholt. Danach ist die
          Wahrscheinlichkeit, da\3 $B$ in Wirklichkeit die Geheimt"ur
          gar nicht "offnen kann und jedesmal nur Gl"uck hatte, nur
          $1/2^n$.
  \end{enumerate}
\end{beispiel}

\begin{definition}[Perfekt zero-knowledge]
Ein Protokoll hei"st {\em perfekt zero-knowledge}, wenn es einen
Simulator gibt, der mit polynomialen Aufwand das Protokoll so
simuliert, da\3 es identisch zur richtigen Protokollausf"uhrung
ist. Der Simulator kennt dabei das Geheiminis nicht.
\end{definition}

Die abstrakte Definition wird weiter unten eingehender erl"autert. Es
l"a\3t sich auch eine etwas schw"achere Definition erstellen, die die
absolute Sicherheit durch {\em nicht berechenbar mit heutigen Mitteln}
ersetzt:

\begin{definition}[Berechenbar zero-knowledge]
Ein Protokoll ist {\em berechenbar zero-knowledge}, wenn es einen
Simulator gibt, der mit polynomialen Aufwand das Protokoll so
simuliert, da\3 es mit heutigen Mitteln nicht von einer richtigen
Protokollausf"uhrug zu unterscheiden ist.
\end{definition}

Berechenbar zero-knowledge bietet zwar nicht den unbedingten Schutz von
perfekt zero-knowledge, da durch neue Algorithmen oder schnellere Computer die
Unterscheidung doch m"oglich werden kann. Es ist in vielen F"allen
das einzige, was "uber ein Protokoll bewiesen werden kann.

Der Simulator ist das zentrale Instrument, mit dem die Zero-knowledge
Eigenschaft definiert und bewiesen wird. Die Aufgabe dieses Simulators
ist es, in Interaktionen mit einem echten Verifizierer die Rolle
des Beweisers zu "ubernehemen. Dabei steht dem Simulator das Geheimnis
des Beweisers nicht zur Verf"ugung. Gefordert ist nun, da\3 er es in
einem nicht vernachl"a\3igbaren (polynomial gro\3en) Teil der
Protokollausf"uhrungen schafft, die erwarteten Antworten zu geben und
das Protokoll damit korrekt ablaufen zu lassen. Der Simulator stellt
einerseits keinen erfolgreichen Angriff auf das Protokoll dar, da er
bei konkreten Implementierungen nur in einem Teil der F"alle die
richtigen Antworten geben kann. Andererseits sind die Antworten, in
denen der Simulator korrekte Antworten geben konnte, nicht von den
Antworten des Beweisers unterscheidbar. Daher kann sich ein
Verifizierer mit Hilfe des Simulators beliebig (polynomial) viele
Ausf"uhrungen des Protokolls beschaffen, ohne mit dem Beweiser
interagieren zu m"ussen. Der Verifizierer erh"alt daher durch den
Ablauf des Protokolls mit dem Beweiser keine Informationen, die er
nicht auch mit dem Simulator h"atte erhalten k"onnen. Deswegen kann
der Beweiser sicher sein, da\3 er keine Information "uber sein
Geheimnis preisgibt, wenn er das Protokoll polynomial oft ausf"uhrt.

\begin{beispiel}
Der Keller mit Geheimt"ur ist perfekt zero-knowledge. Was w"urde
passieren, wenn der Eingang keinen Knick machen w"urde?
\end{beispiel}

Ein Beispiel f"ur einen Zero-knowledge Beweis wird am folgenden
Protokoll, mit dem ein Beweiser nachweist, die L"osung des
Rubik W"urfels aus einer bestimmen Stellung zu kennen, gezeigt.

\begin{beispiel}[Rubik W"urfel]
$B$ will $V$ beweisen, da\3 er den W"urfel aus einer bestimmten
Stellung $P$ l"osen kann. Es wird davon ausgegangen, da\3 es 
schwierig ist, eine allgemeine L"osung des W"urfel zu finden.
\begin{enumerate}
  \item B transformiert den W"urfel in eine andere Stellung $P'$ mit 
        Hilfe einer Transformation $\tau$.
  \item V entscheidet sich, ob er die L"osung aus der Stellung $P'$ 
        heraus oder die Transformation $\tau : P \leftrightarrow P'$ 
        sehen will.
  \item B f"uhrt das Verlangte vor.
\end{enumerate}
Ist der Beweiser echt, so kann er auf jede der beiden m"oglichen
Fragen des Verifiziers eine korrekte Antwort geben. Ist er jedoch ein
Betr"uger, so kann er auf die Aufgabe "`L"ose den W"urfel aus der
Stellung $P'$"' nicht l"osen, da er daf"ur eine allgemeine L"osung des
Rubik W"urfels kennen m"u\3te, was aber nach der Grundannahme des
Protokolls schwierig ist. Die Aufgabe "`Zeige die Transformation 
$\tau : P \leftrightarrow P'$"' kann er trivialerweise l"osen.

Um zu ermitteln, ob das Protokoll zero-Knowledge ist, mu\3 ein
Simulator angegeben werden. 

Der Simulator hat zwei Rubik W"urfel, die sich im Anfangszustand bzw. im 
Zustand $P$ befinden. Er r"at zuerst die Herausforderung des
Verifiziers. Nun kann er das Problem $P'$ so konstruieren, da\3 er die
Antwort kennt, falls sich der Verifizierer f"ur die Herausforderung
entscheidet, die er vorher geraten hat. Nachdem der Simulator nun das
Problem~$P'$ pr"asentiert hat, fordert der Verifizierer den Simulator
eine der L"osungen anzugegen. Stimmt die Herausforderung mit der
geratenen "uberein, so gibt der Simulator die korrekte
Antwort. Stimmen die Herausforderungen nicht "uberein, so gibt der
Simulator auf, und fordert zu einem neuen Protokolldurchlauf auf.

Dieses Protokoll ist zwar perfekt zero-knowledge, {\em aber} der 
Verifizier kann aus dem Weg die L"osung erraten.
\end{beispiel}

Aus diesem Beispiel l"a\3t sich die typische Strukur eines
kryptographisch nutzbaren Zero-knowledge Protokolls ableiten.

{\bf Typische Struktur eines zero-knowledge Protokolls:}

Der Beweiser $B$ will dem Verifizierer $V$ beweisen, die L"osung eines
Problems $P$ aus einer Problemklasse ${\cal P}$ zu kennen:
\begin{enumerate}
  \item B bestimmt ein zu $P$ "aquivalentes Problen $P'\in {\cal P}$
        mit Hilfe einer Transformation $\tau : P \mapsto P'$ und
        sendet $P'$ als Referenzwert an V.
  \item V kann sich f"ur eine Herausforderung (challenge) $r$
        entscheiden. Im Regelfall ist dies die Frage, ob er die
        L"osung des Problems $P'$ oder die Transformation $\tau : P
        \mapsto P'$ sehen will.
  \item B sendet die zum Referenzwert $P'$ und der Herausforderung $r$
        passende Antwort.
\end{enumerate}

Kennt der Beweiser die Herausforderung $r$ bevor er seinen
Referenzwert $P'$ offenlegen mu\3, so braucht er die L"osung gar nicht
zu kennen, da er sich zuerst die Antwort w"ahlen kann und dann einen zur
Herausforderung passende Referenzwert w"ahlen kann. Daraus ergibt sich

{\bf die Strukur eines Simulators:}
\begin{enumerate}
  \item W"ahle die Herausforderung $r'$ zuf"allig und gleichverteilt.
  \item Berechne einen zu $r'$ passenden Referenzwert und eine 
        dazugeh"orige Antwort.
  \item Der Verifizierer stellt die Herausforderung $r$.
  \item Falls $r \neq r'$ gehe zu Schritt 1.
  \item Gebe den Protokollablauf mit Referenzwert, Herausforderung und 
        Antwort aus.
  \item Weiter mit Schritt 1.
\end{enumerate}
Eine M"oglichkeit f"ur die Realisierung von Zero-Knowledge Protokollen
ist relativ offensichtlich. Diese M"oglichkeit basiert auf einem
Beweiser, der behauptet eine L"osung einer Klasse von Problemen zu
besitzen. Der Beweis kann dann einfach angetreten werden, indem der
Verifizierer ein Problem aus dieser Klasse benennt und der Beweiser
eine L"osung f"ur dieses Problem vorlegt.
\begin{beispiel}[Faktorisieren] 
Der Beweiser behauptet, effizient gro\3e Zahlen faktorisieren zu k"onnen.
  \begin{enumerate}
    \item V gibt eine gro\3e Zahl an.
    \item B faktorisiert sie und V "uberpr"uft das Ergebnis.
  \end{enumerate}
Der Simulator S w"ahlt sich einige Primzahlen und berechnet das Produkt
$r'$. Falls der Verifizierer V eine Zahl $r$ w"ahlt mit $r=r'$, so
kann S ihre Faktorisierung angeben. Das Problem ist es, da"s
die Wahrscheinlichkeit f"ur $r=r'$ unendlich klein ist, da V jede
der unendlich vielen Zahlen w"ahlen kann. Gefordert ist aber, da"s der
Simulator in einem nicht vernachl"a"sigbaren Teil der 
Protokollausf"uhrungen die erwarteten Antworten geben kann.
\end{beispiel}

\begin{beispiel}[Kenntnis von Quadratwurzeln]
Es sei $n$ ein \"offentlich bekannter, ge\-n\"u\-gend gro"ser Modulus. 
Die Faktorisierung
sei nicht "offentlich bekannt. Das Berechnen von Wurzeln modulo $n$
ist "aquivalent zu der Faktorisierung von $n$ und damit schwierig. Es
sei $y_B$ ein "offentlich bekannter Wert. Der Beweiser $B$ behauptet, ein
$x_B$ mit $x_B^2 \equiv y_B \bmod n$ zu kennen.  Das folgende
Protokoll kann von ihnen benutzt werden, um diese Behauptung zu
"uberpr"ufen.
\begin{enumerate}
   \item B w"ahlt x zuf"allig und legt $y \equiv x^2 \bmod n$ offen.
   \item V w"ahlt sich $r \in \{ 0,1 \}$ und teilt es $B$ mit.
   \item B berechnet  $z \equiv x_B^r x \bmod n$.
   \item V "uberpr"uft ob $z^2 \equiv y_B^r y \bmod n$.
\end{enumerate}
Durch die Multipikation mit $x$ wird das urspr"ungliche Problem auf
ein neues Problem, die Quadratwurzel aus $y_B y$ zu ziehen,
transformiert. Falls V diese Quadratwurzel sehen will, kann B sie
berechnen. Falls aber V die Transformation sehen will, so liefert V
einfach die Zahl x. Die Transformation ist hier auch durch
Wurzelziehen zu berechen und deshalb genauso schwierig wie das
urspr"ungliche Problem.

Um einen Simulator zu konstruieren, halten wir uns an die oben 
beschriebene typische Struktur eines Simulators. Der Simulator w"ahlt 
zuf"allig und gleichverteilt die Herausforderung $r'\in\left\{0,1\right\}$. 
Falls der Verifizierer die andere Herausforderung $r\neq r'$ w"ahlt, gibt der 
Simulator auf. Sonst kann er die richtigen Antworten geben:

Simulator S: $r'=0$
\begin{enumerate}
  \item S w"ahlt $x$ zuf"allig und legt $y\equiv x^2$ offen.
  \item V entscheidet sich f"ur $r=0$.
  \item S berechnet $z:=x_B^0 x \equiv x$.
  \item V "uberpr"uft, ob $z^2\equiv y_B^0 y$.
\end{enumerate}
Simulator S: $r'=1$
\begin{enumerate}
  \item S w"ahlt $x$ zuf"allig und legt $y\equiv y_B^{-1} x^2$ offen.
  \item V entscheidet sich f"ur $r=1$.
  \item S berechnet $z:=x$.
  \item V "uberpr"uft, ob $z^2\equiv y_B^1 y$.
\end{enumerate}

In der H"alfte der F"alle liefert der Simulator richtige Antworten, d.h.
das Protokoll ist ein Zero-knowledge-Protokoll. Es ist sogar perfekt
zero knowledge, da der Simulator nur quadratische Reste ver"offentlicht.
\end{beispiel}
\begin{beispiel}[Graphenisomorphie]
Der Beweiser kennt einen Isomorphismus $\phi$ zwischen den Graphen
$G_1$ und $G_2$ und will dies dem Verifizierer beweisen. Dazu eignet
sich folgendes Protokoll:
\begin{enumerate}
  \item B transformiert den Graphen $G_1$ durch eine zuf"allige 
        Permutation $\tau$ und erh"alt einen neuen Graphen $H=\tau(G_1)$.
  \item V w"ahlt, ob er $G_1 \mapsto H$ oder $G_2 \mapsto H$ sehen will.
  \item B gibt den geforderten Isomorphismus an, da er die beiden Isomorphismen
        \begin{eqnarray*}
          \tau                  & : & G_1 \longrightarrow H \\
          \tau \circ \phi^{-1}  & : & G_2 \longrightarrow H 
        \end{eqnarray*}
        kennt.
\end{enumerate}
Das Protokoll ist perfekt zero-knowledge.
\end{beispiel}




   
  

